\documentclass[runningheads]{llncs}
\usepackage[T1]{fontenc}
\usepackage{graphicx}
\usepackage{booktabs}
\usepackage{hyperref}

\begin{document}

\title{Scribble-Conformant Cached Segmentation: A Data-Centric Approach to
Interactive Whole-Body PET/CT Lesion Segmentation}

\author{Usi Adia-Nimuwa}
\authorrunning{U. Adia-Nimuwa}
\institute{Independent researcher\\
\email{usi.jadia@gmail.com}\\
Code: \url{https://github.com/il-miscusi/autopetv-algorithm}}

\maketitle

\begin{abstract}
We describe a data-centric submission to the autoPET~V challenge on
interactive, clinician-in-the-loop lesion segmentation in whole-body PET/CT.
The organizers' pretrained four-channel nnU-Net baseline is left unchanged;
all improvement comes from how its predictions are managed across the
interactive loop. The previous iteration's prediction is cached and each new
corrective scribble is enforced by construction: foreground scribbles trigger
PET-adaptive connected region growing united with a locally-masked model
re-run, while background scribbles trigger component removal discriminated by
uptake statistics, with a scribble-commensurate, depth-gated shell chip
protecting mostly-true components. Because the evaluation places each
scribble on the largest error component of the previous prediction, every
iteration reduces its target error and the per-case Dice trajectory is
approximately monotone, in contrast to the stateless baseline, whose
trajectory regresses between iterations. The approach requires no training,
runs in seconds for correction-only iterations, and conforms to the
challenge's data-centric award track.
\keywords{Interactive segmentation \and PET/CT \and Lesion segmentation
\and Scribbles \and Data-centric AI.}
\end{abstract}

\section{Introduction}

autoPET~V evaluates algorithms that produce an initial whole-body PET/CT
lesion segmentation and refine it over five corrective steps, where each
step supplies one simulated or clinician-drawn scribble placed inside the
largest error region of the previous prediction: a foreground (tumor)
scribble when the largest miss outweighs the largest false positive, and a
background scribble otherwise. Cases are scored by the area under the
per-iteration Dice and lesion-detection curves, so per-scribble
responsiveness is worth as much as raw single-shot quality.

The organizers' baseline concatenates Gaussian scribble heatmaps as third
and fourth input channels to an nnU-Net~\cite{isensee2021nnu} trained on the
FDG~\cite{gatidis2022whole} and PSMA~\cite{jeblick2024psma} databases with
simulated scribbles. At test time it re-runs inference from scratch at every
iteration. Two properties limit it: the network may ignore the scribble
channels (nothing enforces conformance), and statelessness discards the
previous prediction, so Dice can regress between iterations---the baseline's
own reference trajectory on its sample case is
$0.870 \rightarrow 0.886 \rightarrow 0.871 \rightarrow 0.878 \rightarrow
0.888$.

Our submission treats these as the whole problem. We change no weights and
train nothing; we manage predictions across the loop.

\section{Method}

\subsection{Cached interactive state}

The Grand Challenge interface invokes the container once per iteration and
provides a per-case scratch directory. We cache the previous binary
prediction and the previous clicks JSON there. Scribbles accumulate
append-only per class, so the new scribble is recovered as the per-class
suffix of the current clicks relative to the cached clicks. Iteration~0
(no clicks) runs the unchanged baseline nnU-Net with zeroed heatmap
channels. On a cold cache with clicks present (defensive path), the baseline
is re-run with heatmaps and all scribbles are conformed at once.

\subsection{Foreground scribbles: grow, then locally trust the model}

A foreground scribble traces the centerline (or boundary/random subset) of
the largest cross-section of the largest missed lesion region, and every
scribble voxel is ground-truth tumor. We union into the cached prediction:
(i)~a connected region grown from the scribble in the PET volume, thresholded
at an escalating fraction (0.45--0.8) of the median seed uptake within a
padded local box, accepted at the first threshold whose region stays under
400\,ml, falling back to a small dilation of the scribble; and (ii)~the
voxels that a re-run of the baseline (with the accumulated scribble
heatmaps) adds within 60\,mm of the new scribble --- the model contributes
the lesion's three-dimensional shape where thresholding alone is crude,
but only locally, so it can never retract or alter the prediction elsewhere.
Model re-runs are skipped when the invocation approaches the per-job time
limit; the growth step alone still guarantees progress.

\subsection{Background scribbles: discriminate before deleting}

A background scribble lies in the largest false-positive region, but the
predicted component containing it is not necessarily wholly false: rim-like
oversegmentation of a correctly detected lesion mass places the scribble
inside a component that is mostly true, and deleting that component is
catastrophic (measured Dice $0.87 \rightarrow 0.10$ on the organizers'
sample case). We therefore remove the whole component only when it is small
($<$30\,ml) or when, with PET available, the scribble's median uptake reaches
half of the component's 75th-percentile uptake --- uniformly hot
physiological structures (bladder, kidneys, brain) satisfy this, whereas a
scribble in a cold oversegmentation rim around a hot true lesion does not.
Otherwise we chip conservatively: an ellipsoidal neighbourhood of the
scribble (4\,mm in-plane, 6\,mm axially), further gated to voxels whose
depth inside the component does not exceed the scribble's own depth plus
4\,mm. The removal is thereby commensurate with the scribble --- a
15-voxel rim sliver removes a shell, not a 25\,mm ball (which measurably
deleted 2{,}013 true voxels to eliminate 229 false ones).

Finally, all foreground scribble voxels are set positive and all background
scribble voxels negative, unconditionally --- they are the only voxels whose
labels are certain.

\section{Results}

\paragraph{Sample-case validation.} Replaying the organizers' evaluator on
their released sample case and baseline outputs (foreground growth and model
re-runs disabled --- a lower bound using only removal, fallback dilation and
enforcement), our trajectories are approximately monotone with Dice AUC
3.484--3.487 across the three scribble strategies versus 3.514 for the full
baseline; naive component deletion, by contrast, collapses to AUC 1.65.
% TODO: DeepPSMA full-pipeline table (baseline vs ours, centerline/random)
% when the GPU validation run completes.

\paragraph{Platform validation.} On the preliminary (non-interactive,
single-shot) test phase, the submission reproduces baseline behaviour as
expected (mean Dice 0.647, mean detection F1 0.730 over five cases),
confirming interface and runtime conformance: iteration~0 is one nnU-Net
inference, correction-only iterations complete in seconds, and the per-job
limit is never approached.

\section{Discussion}

The evaluation protocol makes each scribble an oracle annotation of the
previous prediction's largest error. A stateless model may or may not act on
it; a cached, conformant post-processor acts on it by construction, and its
gains compound across the loop because the corrected error does not
reappear. The design errs toward locality: it never lets new inference
retract distant regions, accepting slower global improvement in exchange for
monotone trajectories, which is the correct trade under an AUC-over-
iterations metric. All improvement is pre/post-processing around the
unchanged baseline model, per the challenge's data-centric track.

\begin{thebibliography}{4}
\bibitem{isensee2021nnu}
Isensee, F., Jaeger, P.F., Kohl, S.A.A., Petersen, J., Maier-Hein, K.H.:
nnU-Net: a self-configuring method for deep learning-based biomedical image
segmentation. Nature Methods 18(2), 203--211 (2021)

\bibitem{gatidis2022whole}
Gatidis, S., et al.: A whole-body FDG-PET/CT dataset with manually annotated
tumor lesions. Scientific Data 9, 601 (2022)

\bibitem{jeblick2024psma}
Jeblick, K., et al.: A whole-body PSMA-PET/CT dataset with manually
annotated tumor lesions. The Cancer Imaging Archive (2024)

\bibitem{autopetv2026}
autoPET~V challenge: Automated Lesion Segmentation in Whole-Body PET/CT ---
the human-AI frontier. \url{https://autopet-v.grand-challenge.org} (2026)
\end{thebibliography}

\end{document}
